```latex
\documentclass[11pt,a4paper]{article}

\usepackage{amsmath,amssymb,amsthm}
\usepackage[margin=1.1in]{geometry}
\usepackage{hyperref}
\usepackage{booktabs}

\hypersetup{colorlinks=true,linkcolor=blue,citecolor=blue,urlcolor=blue}

\theoremstyle{plain}
\newtheorem{theorem}{Theorem}[section]
\newtheorem{proposition}[theorem]{Proposition}
\newtheorem{lemma}[theorem]{Lemma}
\newtheorem{corollary}[theorem]{Corollary}
\theoremstyle{definition}
\newtheorem{definition}[theorem]{Definition}
\newtheorem{example}[theorem]{Example}
\theoremstyle{remark}
\newtheorem{remark}[theorem]{Remark}

\newcommand{\Dset}{\mathcal{D}}
\newcommand{\Z}{\mathbb{Z}}
\newcommand{\N}{\mathbb{N}}

\title{A sharp stability threshold for Erd\H{o}s Problem 848:\\
admissible sets meeting both extremal classes\\
have density at most $1/50$}
\author{Paul J.\ Lee}
\date{November 2026}

\begin{document}
\maketitle

\begin{abstract}
Call $A\subseteq\{1,\dots,N\}$ \emph{admissible} if $ab+1$ is not squarefree for all $a,b\in A$,
including $a=b$. A.~C.~Li has determined the maximum size $f(N)=\lfloor (N+18)/25\rfloor$ of an
admissible set for every $N$, in a Lean-formalized proof which we take as an input; the extremal sets
are the residue classes $n\equiv 7 \pmod{25}$ and $n\equiv 18\pmod{25}$. Sawhney had previously proved a
stability statement: there is an (unspecified) $\eta>0$, implicitly $\eta\approx 0.0023$, such that for
large $N$ every admissible set of size at least $(1/25-\eta)N$ lies inside one of these two classes;
he left open the structure of admissible sets of smaller density. We close that gap in sharp form.
We prove that for every $\varepsilon>0$ and all $N\ge N_0(\varepsilon)$, every admissible
$A\subseteq\{1,\dots,N\}$ which is \emph{not} contained in a single one of the two classes satisfies
$|A|\le (1/50+\varepsilon)N$, and that the constant $1/50$ is sharp, attained exactly by the two
residue-class constructions $\{n:\ n\equiv 7,93 \pmod{100}\}$ and $\{n:\ n\equiv 43,57\pmod{100}\}$.
The proof combines uniform analytic density bounds for the relevant squarefree-sieve events with two
explicit finite computations, both described in the paper: the determination of the divisibility
(``gluing'') relation on residues modulo $900$, and the maximization of the resulting weighted
objective by $0$--$1$ programming together with an exhaustive enumeration over the $64$ Galois-closed
signature sets. No explicit $N_0$ is claimed.
\end{abstract}

\section{Introduction}

\subsection{The problem}

For $N\in\N$ write $[N]=\{1,2,\dots,N\}$.

\begin{definition}\label{def:adm}
A set $A\subseteq[N]$ is \emph{admissible} if $ab+1$ is not squarefree for all $a,b\in A$,
\emph{including the diagonal case $a=b$}. We write
\[
 f(N)=\max\{|A|:\ A\subseteq[N]\ \text{admissible}\}.
\]
\end{definition}

Erd\H{o}s Problem 848 (Erd\H{o}s--S\'ark\"ozy) asks whether $f(N)$ is attained by the residue class
$\{n\le N:\ n\equiv 7\pmod{25}\}$; see \cite{EP848}. The diagonal convention $a=b$ is the one used in
Sawhney's proof \cite{Sawhney} and in Li's formalization \cite{Li}; it is essential below, since our
argument extracts information from the non-squarefreeness of $a^2+1$ and gives nothing without it.

Two residue classes are admissible for trivial local reasons:
\[
 \mathcal{C}^{7}_N=\{n\le N:\ n\equiv 7\ (25)\},\qquad
 \mathcal{C}^{18}_N=\{n\le N:\ n\equiv 18\ (25)\},
\]
because $7\cdot 7+1=50$, $18\cdot 18+1=325$ and $7\cdot 18+1=127$ are such that $25\mid ab+1$ whenever
$a\equiv b\equiv 7\ (25)$ or $a\equiv b\equiv 18\ (25)$. (Note that $25\nmid 127$: the two classes are
\emph{not} glued to each other by $25$. This will matter.) We call these the two
\emph{extremal classes}.

\subsection{Inputs: Li's theorem and Sawhney's stability theorem}

\begin{theorem}[A.~C.~Li \cite{Li}]\label{thm:li}
For every $N\ge 1$,
\[
 f(N)=\bigl|\{n\le N:\ n\equiv 7\ (\mathrm{mod}\ 25)\}\bigr|=\Bigl\lfloor \frac{N+18}{25}\Bigr\rfloor .
\]
\end{theorem}

Theorem~\ref{thm:li} is proved in \cite{Li} by an exact prefix colouring for $N\le 5\cdot 10^{6}$
combined with a Hall-defect sieve argument for larger $N$, and is formalized in Lean~4 in the
accompanying repository. It is an \emph{input} to the present note: we neither reprove nor modify it.
(Section~\ref{sec:verif} records an independent source rebuild and external kernel replay of that
formalization, performed by the present author; that work verifies Li's theorem and is logically
separate from the theorem proved here.)

\begin{theorem}[Sawhney \cite{Sawhney}]\label{thm:sawhney}
There is $\eta>0$ such that for all sufficiently large $N$, every admissible $A\subseteq[N]$ with
$|A|\ge (1/25-\eta)N$ satisfies $A\subseteq \mathcal{C}^{7}_N$ or $A\subseteq \mathcal{C}^{18}_N$.
\end{theorem}

Sawhney does not specify $\eta$; his four case bounds are
$0.0377N$, $0.0358N$, $0.0336N$, $0.0294N$ (each $+o(N)$), so that $\eta\approx 0.0023$ is implicit.
He explicitly leaves open ``a strong structural classification of sets $A$ with size $\delta N$''.
Thus before the present note, the structure of admissible sets was understood only in the top
sliver of densities, $|A|/N\in[1/25-0.0023,\ 1/25]$, and nothing was known in the whole middle range
below it.

\subsection{Main theorem}

\begin{theorem}\label{thm:main}
For every $\varepsilon>0$ there is $N_0(\varepsilon)$ such that the following holds for all
$N\ge N_0(\varepsilon)$. If $A\subseteq[N]$ is admissible and $A$ is \emph{not} contained in
$\mathcal{C}^{7}_N$ and \emph{not} contained in $\mathcal{C}^{18}_N$, then
\[
 |A|\ \le\ \Bigl(\frac{1}{50}+\varepsilon\Bigr)N .
\]
Equivalently: every admissible set with more than $(1/50+\varepsilon)N$ elements lies inside one
of the two extremal classes. The constant $1/50$ is sharp: it is attained, within the mod-$900$
type framework of Section~\ref{sec:opt} and up to $O(1)$ errors, by exactly two residue-class
constructions,
\[
 A_1=\{n\le N:\ n\equiv 7\ \text{or}\ 93\ (\mathrm{mod}\ 100)\},\qquad
 A_2=\{n\le N:\ n\equiv 43\ \text{or}\ 57\ (\mathrm{mod}\ 100)\},
\]
both of which are admissible, meet both extremal classes, and have $|A_i|=N/50+O(1)$.
\end{theorem}

Combined with Theorem~\ref{thm:li}, the picture for large $N$ is therefore complete up to
$\varepsilon$: an admissible set either sits inside one extremal class (where its size is at most
$\lfloor (N+18)/25\rfloor$, with equality for the class itself), or it has density at most
$1/50+\varepsilon$, and the latter bound is attained.

\begin{remark}[Ineffectivity]\label{rem:ineff}
No explicit $N_0(\varepsilon)$ is claimed anywhere in this note. The density bounds of
Section~\ref{sec:dens} are stated in an $\varepsilon$--$T$--$N_0$ form in which the truncation
parameter $T=T(\varepsilon)$ and then $N_0$ are chosen but not tracked, and the large-prime error
terms are handled exactly as in Sawhney's argument \cite{Sawhney}, which we take as given and do not
re-derive. In this sense $N_0(\varepsilon)$ is ineffective as stated here, and it inherits the
ineffectivity of the inputs \cite{Li,Sawhney} in the surrounding discussion.
% TODO(author): an effective N_0 would require tracking T(eps), the tail sum over primes q > T, and
% the O_T(1) CRT error; none of the source material does this.
\end{remark}

\begin{remark}[Prior work and what is new]\label{rem:prior}
The observation that the residues $7$ and $93$ modulo $100$ are compatible ``through $4$'', because
$7\cdot 93+1=652=4\cdot 163$, is due to \textbf{AstroCat93} \cite{AstroCat}, whose verification write-up
also isolates a sub-case of limiting density $0.02$; its overall bound is $\approx 0.0282$. That work
does not show that density $1/50$ is \emph{attained} by an admissible set, and does not establish
$1/50$ as the global stability threshold. What is new here is (i) the proof that the full
configuration $\{7,93\}\pmod{100}$ (and its mirror $\{43,57\}\pmod{100}$) is admissible and attains
density exactly $1/50$, and (ii) the matching sharp upper bound of Theorem~\ref{thm:main}, obtained by
reducing all mixed configurations to the single numerical inequality $72K_2'+828J'<18$ of
Section~\ref{sec:opt}. Other community bounds in circulation are weaker: $\approx 0.0276$
(bzanghi/erdos-848) and an existential $\eta=1/2000$ in the Lean work of
The-Obstacle-Is-The-Way/erdos-banger. A novelty search carried out on 27 October 2026 found neither
the $1/50$ threshold nor the two constructions stated anywhere; a few sources (notably
\S IV.1 of arXiv:2511.16072, and some raw forum material) were not read in full.
\end{remark}

\section{The extremal configurations}\label{sec:extremal}

\begin{proposition}\label{prop:extremal}
The sets $A_1=\{n\le N:\ n\equiv 7,93\ (100)\}$ and $A_2=\{n\le N:\ n\equiv 43,57\ (100)\}$ are
admissible, each meets both extremal classes, and $|A_i|=N/50+O(1)$.
\end{proposition}

\begin{proof}
Diagonal and within-residue pairs. If $a\equiv b\equiv 7\ (100)$ then $ab+1\equiv 50\ (100)$, so
$25\mid ab+1$; in particular $7^2+1=50$. If $a\equiv b\equiv 93\ (100)$ then again
$ab+1\equiv 93\cdot 93+1=8650\equiv 50\ (100)$, and $8650=2\cdot 5^2\cdot 173$, so $25\mid ab+1$.
Cross pairs. If $a\equiv 7$ and $b\equiv 93\ (100)$ then $ab+1\equiv 7\cdot 93+1=652\equiv 52\ (100)$,
so $4\mid ab+1$ (indeed $652=2^2\cdot 163$). In all cases $ab+1$ is divisible by $25$ or by $4$, hence
not squarefree, so $A_1$ is admissible.

For $A_2$: $43^2+1=1850=2\cdot 5^2\cdot 37$ and $57^2+1=3250=2\cdot 5^3\cdot 13$ are both divisible by
$25$, so within-residue pairs give $ab+1\equiv 50\ (100)$; and $43\cdot 57+1=2452\equiv 52\ (100)$ is
divisible by $4$. So $A_2$ is admissible.

Finally $7\equiv 7$ and $93\equiv 18\ (25)$, while $43\equiv 18$ and $57\equiv 7\ (25)$; so each $A_i$
meets both extremal classes and is contained in neither. Each $A_i$ consists of two residue classes
mod $100$, whence $|A_i|=2N/100+O(1)=N/50+O(1)$.
\end{proof}

\begin{proposition}[Finite check]\label{prop:only-two}
Among all pairs $\{u,v\}$ of residues mod $100$ with $u\equiv 7$, $v\equiv 18\pmod{25}$ (or vice
versa), the only pairs for which $\{n\le N:\ n\equiv u\ \text{or}\ v\ (100)\}$ is admissible are
$\{7,93\}$ and $\{43,57\}$.
\end{proposition}

\begin{proof}
There are $4\cdot 4=16$ such pairs, and for each the three conditions ($u^2+1$, $v^2+1$, $uv+1$ not
squarefree by a local obstruction at $4$, $9$ or $25$) are decided by a direct computation modulo
$900$. The enumeration was carried out by the author and independently by the referee of
\cite{EP848Code}, and also checked against all actual integer pairs up to $20000$.
\end{proof}

\section{Types modulo $900$ and gluing}\label{sec:types}

Throughout, set
\[
 M=900=4\cdot 9\cdot 25 .
\]
The \emph{type} of an integer $n$ is its residue $n\bmod M$. Let
\[
 \Dset=\{p\ \text{prime}:\ p\equiv 1\ (\mathrm{mod}\ 4),\ p\ge 13\}.
\]
For integers $a,b$ define the two sieve events
\[
 D(a):\ \exists p\in\Dset,\ p^2\mid a^2+1;\qquad
 P_b(a):\ \exists q\ \text{prime},\ q\ge 7,\ q^2\mid ab+1 .
\]

\begin{lemma}[Local facts; finite verification]\label{lem:F1}
\begin{enumerate}
\item[(F1)] For every $n$, $n^2+1$ is divisible by neither $4$ nor $9$; it is divisible by $25$ if and
only if $n\equiv 7$ or $18\pmod{25}$; no prime $p\equiv 3\pmod 4$ divides $n^2+1$; and for
$p\equiv 1\pmod 4$ the congruence $x^2\equiv -1\pmod{p^2}$ has exactly two roots. Consequently, if
$n\bmod 25\notin\{7,18\}$, then
\[
 n^2+1\ \text{is not squarefree}\iff D(n).
\]
\item[(F2)] Whether $4$, $9$ or $25$ divides $ab+1$ depends only on the types of $a$ and $b$. Call a
pair of types $(s,t)$ \emph{glued} if $4\mid st+1$ or $9\mid st+1$ or $25\mid st+1$. If $(s,t)$ is not
glued, then for $a$ of type $s$ and $b$ of type $t$,
\[
 ab+1\ \text{is not squarefree}\iff P_b(a),
\]
and necessarily $q\nmid b$ for the witnessing prime $q$.
\item[(F3)] Let $C\subseteq\Z/M\Z$ be the set of $72$ types $s$ with $s\equiv 7$ or $18\pmod{25}$.
Every $s\in C$ is glued to itself and to every type in the same class mod $25$. A type outside $C$ is
never glued to itself. A type in $C$ and a type outside $C$ can be glued only through $4$ or $9$;
likewise a $7$-class type and an $18$-class type can be glued only through $4$ or $9$.
\end{enumerate}
\end{lemma}

\begin{proof}
(F1) Squares mod $4$ lie in $\{0,1\}$ and squares mod $9$ in $\{0,1,4,7\}$, so $n^2+1\not\equiv 0$
mod $4$ or $9$. The roots of $x^2\equiv-1\pmod{25}$ are $x\equiv 7,18$. If $p\equiv 3\pmod 4$ then
$-1$ is a non-residue mod $p$. If $p\equiv 1\pmod 4$, Hensel lifting gives exactly two roots mod
$p^2$. Hence, for $n$ with $n\bmod 25\notin\{7,18\}$, the only primes whose square can divide $n^2+1$
are $p\equiv 1\pmod 4$ with $p\ne 5$, i.e.\ $p\in\Dset$.

(F2) The first assertion is clear since $4,9,25\mid M$. If $(s,t)$ is not glued, then $4,9,25\nmid
ab+1$; since $2^2=4$, $3^2=9$, $5^2=25$, any prime square dividing $ab+1$ has prime $q\ge 7$. If
$q\mid b$ and $q^2\mid ab+1$ then $q\mid 1$, a contradiction.

(F3) If $s\equiv t\equiv 7\ (25)$ then $st+1\equiv 50\equiv 0\ (25)$, and similarly
$18\cdot 18+1=325\equiv 0\ (25)$; this covers self-gluing and same-class gluing. If $s\notin C$ then
$s^2+1$ is divisible by none of $4,9,25$ by (F1). Finally, if $25\mid st+1$ with $s\equiv 7\ (25)$ then
$t\equiv-7^{-1}\equiv-18\equiv 7\ (25)$ (as $7\cdot 18=126\equiv 1$), so $t$ is in the same class; the
same computation with $7$ and $18$ interchanged gives the remaining statement.
\end{proof}

The local facts were verified computationally as well: (F1) for all primes $p<2000$, (F2) on $20{,}000$
random pairs, and (F3) on all $900\times 900$ type pairs (see Section~\ref{sec:verif}).

\section{Uniform density bounds}\label{sec:dens}

All bounds in this section are uniform in the ``pivot'' $b\le N$, which is what allows us to choose
the witness $b$ depending on the set $A$.

\begin{lemma}[Truncation]\label{lem:trunc}
Let $T\ge 7$ and let $E$ be either of the events $P_b$, $D\wedge P_b$. For a fixed residue $s$ mod $M$,
the count $\#\{a\le N:\ a\equiv s\ (M),\ E(a)\}$ differs from the corresponding count with all
witnessing primes restricted to $q\le T$ by at most
\[
 2\sum_{T<q\le N}\Bigl(\frac{N}{q^2}+1\Bigr)\ \le\ \frac{2N}{T}+2\pi(N),
\]
and the truncated count is an exact CRT count $\mathrm{dens}_T\cdot N/M+O_T(1)$. Only primes $q\le N$
need be considered, since $q^2\le ab+1\le N^2+1$.
\end{lemma}

Given $\varepsilon>0$ we first choose $T=T(\varepsilon)$ so that the truncated CRT densities are within
$\varepsilon/2$ of their limits and $2/T<\varepsilon/4$, and then $N_0$ so that
$2\pi(N)+O_T(1)<\varepsilon N/4$ for $N\ge N_0$. All constants below are limits of truncated products,
and only upper bounds on them are used.

\begin{lemma}[The $K_2'$ cap]\label{lem:K2}
For every $\varepsilon>0$ there is $N_0$ such that for all $N\ge N_0$, all $b\le N$ and all residues
$s$ mod $M$,
\[
 \#\{a\le N:\ a\equiv s\ (M),\ P_b(a)\}\ \le\ (K_2'+\varepsilon)\frac{N}{M},
 \qquad
 K_2'=1-\prod_{q\ge 7}\bigl(1-q^{-2}\bigr)=1-\frac{9.375}{\pi^2}\approx 0.0501139 .
\]
\end{lemma}

\begin{proof}
By (F2) a witnessing prime satisfies $q\ge 7$ and $q\nmid b$; for such $q$ the condition
$q^2\mid ab+1$ puts $a$ in a single residue class mod $q^2$, and the classes for distinct $q$ are
independent by CRT. Hence the truncated density of $P_b$ is
$1-\prod_{7\le q\le T,\,q\nmid b}(1-q^{-2})\le 1-\prod_{7\le q\le T}(1-q^{-2})$, and letting
$T\to\infty$ in the majorant gives $K_2'$. Finally
$\prod_{q\ge 7}(1-q^{-2})=\bigl(6/\pi^2\bigr)\big/\bigl(\tfrac34\cdot\tfrac89\cdot\tfrac{24}{25}\bigr)
=9.375/\pi^2$, since $\tfrac34\cdot\tfrac89\cdot\tfrac{24}{25}=0.64$.\footnote{%
This closed form, $K_2'=1-9.375/\pi^2$, is the correct one; an expanded product form containing a
spurious factor $15/16$ appeared in an early draft of this note and should be disregarded.}
Apply Lemma~\ref{lem:trunc}.
\end{proof}

\begin{lemma}[The $J'$ cap]\label{lem:J}
Put
\[
 A_{\Dset}=\prod_{p\in\Dset}\bigl(1-2p^{-2}\bigr),\quad
 B_{\Dset}=\prod_{p\in\Dset}\bigl(1-p^{-2}\bigr),\quad
 Q=\prod_{\substack{p\equiv 3\,(4)\\ p\ge 7}}\bigl(1-p^{-2}\bigr),
\]
\[
 K_1=1-A_{\Dset}\approx 0.0273465,\qquad
 J'=K_1-Q\,\bigl(B_{\Dset}-A_{\Dset}\bigr)\approx 0.0142450 .
\]
Then for every $\varepsilon>0$ there is $N_0$ such that for all $N\ge N_0$, all $b\le N$ and all
residues $s$ mod $M$ with $s\notin C$,
\[
 \#\{a\le N:\ a\equiv s\ (M),\ D(a)\wedge P_b(a)\}\ \le\ (J'+\varepsilon)\frac{N}{M}.
\]
\end{lemma}

\begin{proof}
By inclusion--exclusion, for the truncated events,
\[
 \mathrm{dens}\bigl(D\wedge P_b\bigr)=1-\prod_\ell Y_\ell-\prod_\ell X_\ell+\prod_\ell Z_\ell,
\]
where $X_\ell, Y_\ell, Z_\ell$ are the local densities at $\ell$ of $\neg P_b$, $\neg D$ and
$\neg D\wedge\neg P_b$ respectively, and $Z_\ell\le X_\ell$ at every prime. Explicitly: for
$\ell\in\Dset$ the set $D_\ell=\{\pm r\}$ with $r^2\equiv-1\ (\ell^2)$ has two elements mod $\ell^2$,
and $P_\ell=\{-b^{-1}\bmod \ell^2\}$ is a single class when $\ell\nmid b$, empty when $\ell\mid b$; for
$\ell\ge 7$, $\ell\notin\Dset$ (i.e.\ $\ell\equiv 3\ (4)$, $\ell\ge 7$) the event $D$ is locally empty.

The density is a product of positive local factors, so the dependence on $b$ may be analysed one prime
at a time, all other factors held fixed, and the signs of the two possible changes do not depend on the
other factors:
\begin{itemize}
\item If $\ell\mid b$ then $P_\ell=\emptyset$, and the change in the density relative to
$\ell\nmid b$ is $(Z'-X')/\ell^2\le 0$ if $P_\ell$ would lie outside $D_\ell$, and $-X'/\ell^2\le 0$ if
it would lie inside, where $X',Z'$ denote the products of the remaining local factors. So
$\gcd(b,\ell)=1$ is the worse (larger) case.
\item If $\ell\in\Dset$ and $P_\ell\subset D_\ell$ (equivalently $\ell^2\mid b^2+1$), then $Z_\ell$
rises from $1-3/\ell^{2}$ to $1-2/\ell^{2}$, which raises the density. So coincidence is the worse case.
\end{itemize}
Taking the worst case at every prime simultaneously therefore majorizes every $b$. In that worst case
$\prod Y=A_{\Dset}$, $\prod X=Q\,B_{\Dset}=1-K_2'$ and $\prod Z=Q\,A_{\Dset}$, giving
\[
 1-A_{\Dset}-Q B_{\Dset}+Q A_{\Dset}=K_1-Q\bigl(B_{\Dset}-A_{\Dset}\bigr)=J' .
\]
Finally $s\notin C$ guarantees, by (F1), that non-squarefreeness of $a^2+1$ is exactly the event $D(a)$,
so the bound is applicable. Apply Lemma~\ref{lem:trunc}.
% TODO(author): the per-prime sign computation above, together with the finite CRT-period checks
% reported below, is exactly what the source material establishes; a fully written-out global
% monotonicity induction over all primes is not present in the sources and should be supplied in a
% final version.
\end{proof}

The product formula underlying Lemma~\ref{lem:J} was checked exactly against full CRT-period counts on
the prime sets $\{7,13\}$, $\{7,13,17\}$, $\{11,13,17\}$ and $\{7,11,13\}$, in each case with the
maximum over $b$ equal to the predicted worst case. Empirically, for $a\le 10^{7}$ and pivots
$b\in\{1,\,30030,\,9699690,\,9999991,\,99,\,3141,\,33459226\}$ (exhibiting $0$ to $3$ coincidences
$\ell^2\mid b^2+1$), all observed densities matched the exact per-$b$ predictions ($|z|\le 0.4$), the
largest observed $D\wedge P_b$ density being $0.01120<J'$. Pivots with $\ge 4$ coincidences were not
tested.

\begin{table}[ht]
\centering
\begin{tabular}{llll}
\toprule
Constant & Definition & Value / enclosure \\
\midrule
$A_{\Dset}$ & $\prod_{p\in\Dset}(1-2p^{-2})$ & $0.9726535060$\\
$B_{\Dset}$ & $\prod_{p\in\Dset}(1-p^{-2})$ & $0.9862566747$\\
$Q$ & $\prod_{p\equiv 3(4),\,p\ge 7}(1-p^{-2})$ & $0.9631226050$\\
$K_1$ & $1-A_{\Dset}$ & $\approx 0.0273465$, \ $[0.02734649,\,0.02734654]$\\
$K_2'$ & $1-\prod_{q\ge 7}(1-q^{-2})=1-9.375/\pi^2$ & $\approx 0.0501139$, \ $[0.05011390,\,0.05011395]$\\
$J'$ & $K_1-Q(B_{\Dset}-A_{\Dset})$ & $\approx 0.0142450$, \ $[0.01424497,\,0.01424500]$\\
generic $b$ & density of $D\wedge P_b$, no coincidences & $\approx 0.0012774$\\
\bottomrule
\end{tabular}
\caption{The density constants. Enclosures were obtained from partial products over primes up to
$P=5\cdot 10^{7}$ with rigorous tails (using $\sum_{n>P}n^{-2}<1/(P-1)$); the author's independent
computation used $P=2\cdot 10^{7}$ and agrees to the printed digits.}
\label{tab:consts}
\end{table}

\section{The per-type bound}\label{sec:pertype}

Let $A\subseteq[N]$ be admissible and let $S\subseteq\Z/M\Z$ be the set of types occurring in $A$. For
$s\in S$ define the weight
\[
 f(s,S)=
 \begin{cases}
 1, & s\in C \text{ and } s \text{ is glued to every } t\in S,\\[2pt]
 K_2', & s\in C \text{ and some } t\in S \text{ is not glued to } s,\\[2pt]
 J', & s\notin C .
 \end{cases}
\]

\begin{proposition}\label{prop:pertype}
For every $\varepsilon>0$ there is $N_0(\varepsilon)$ such that for all $N\ge N_0$ and all admissible
$A\subseteq[N]$ with type set $S$,
\[
 |A|\ \le\ \frac{N}{M}\sum_{s\in S}f(s,S)\ +\ \varepsilon N\ +\ O(M).
\]
\end{proposition}

\begin{proof}
Write $A_s=\{a\in A:\ a\equiv s\ (M)\}$, so $|A|=\sum_{s\in S}|A_s|$. There are at most $M$ summands,
so it suffices to bound each $|A_s|$ by $f(s,S)N/M+\varepsilon N/M+O(1)$.

If $s\in C$ is glued to every $t\in S$ we use the trivial bound $|A_s|\le N/M+1$.

If $s\in C$ and some $t\in S$ is not glued to $s$, pick $b\in A$ of type $t$. By (F3) a $C$-type is
glued to itself, so $t\ne s$ and hence $b\notin A_s$. For every $a\in A_s$, $ab+1$ is not squarefree
and $(s,t)$ is not glued, so $P_b(a)$ holds by (F2). Lemma~\ref{lem:K2} gives
$|A_s|\le (K_2'+\varepsilon)N/M$.

If $s\notin C$ and $|A_s|\ge 2$, pick $b\in A_s$. By (F3), $(s,s)$ is not glued, so every
$a\in A_s\setminus\{b\}$ satisfies $P_b(a)$ by (F2); and $a^2+1$ is not squarefree (diagonal
admissibility) with $s\notin C$, so $D(a)$ holds by (F1). Lemma~\ref{lem:J} gives
$|A_s|\le 1+(J'+\varepsilon)N/M$. If $|A_s|\le 1$ its contribution is at most $1$, which is absorbed in
the $O(M)$ term.
\end{proof}

\section{The maximization}\label{sec:opt}

Set
\[
 F(S)=\frac{1}{M}\sum_{s\in S}f(s,S),
\]
and let $\Sigma$ be the family of all $S\subseteq\Z/M\Z$ that are not contained in the $7$-class and
not contained in the $18$-class (i.e.\ $\{s\in S: s\not\equiv 7\,(25)\}\ne\emptyset$ and
$\{s\in S: s\not\equiv 18\,(25)\}\ne\emptyset$). We must compute $\max_{S\in\Sigma}F(S)$.

\subsection{Signature reduction}

For $s\in\Z/M\Z$ define the \emph{signature} $\sigma(s)=(s\bmod 4,\ s\bmod 9)\in\Z/4\times\Z/9$; there
are $36$ signatures, each containing $25$ types, of which $2$ lie in $C$ and $23$ do not. Thus $C$ has
$72$ elements, parametrized by (class $\in\{7,18\}$, signature), and the complement of $C$ consists of
$36$ groups of $23$ types.

By (F3), gluing within a mod-$25$ class is automatic, while gluing between the two classes, and between
a $C$-type and a non-$C$ type, can only occur through $4$ or $9$ and therefore depends only on
signatures: define the symmetric relation
\[
 (a,b)\ R\ (a',b')\quad\Longleftrightarrow\quad 4\mid aa'+1\ \ \text{or}\ \ 9\mid bb'+1 .
\]
That $R$ computes gluing correctly in all the relevant cases was verified exhaustively over all
$900\times 900$ type pairs. Moreover all $23$ non-$C$ types of a given signature behave identically
towards every $C$-type (verified exhaustively), and each carries the same weight $J'$; so a
configuration is determined, for the purposes of $F$, by
\[
 (G_7,\ G_{18},\ Y)\subseteq (\Z/4\times\Z/9)^3,
\]
the signature sets of the $7$-class types present, of the $18$-class types present, and of the non-$C$
groups present. For a signature set $X$ write $R(X)=\bigcap_{i\in X}R(i)$ (with $R(\emptyset)$ the full
set of $36$ signatures). A signature set is \emph{Galois-closed} if it is of the form $R(X)$,
equivalently $R(R(X))=X$.

\begin{lemma}[Closure reduction]\label{lem:closure}
The maximum of $F$ over $\Sigma$ is attained at a configuration whose governing signature sets are
Galois-closed. There are exactly $64$ Galois-closed signature sets, and exactly $439$ pairs of them
satisfying the compatibility constraints.
\end{lemma}

% TODO(author): the closure reduction is stated by the independent referee and implemented in
% opt_closure.py (which verifies R(R(X)) = X for each enumerated closed set and cross-validates
% F against its definition on all 900 types, including 3000 random type sets); a written-out
% proof of the reduction is not present in the source material and should be supplied in a final
% version. The computational conclusions of Theorem 6.2 below were obtained by three independent
% solvers, only one of which uses the reduction.

\subsection{The computation and its outcome}

\begin{theorem}[Computational]\label{thm:opt}
\[
 \max_{S\in\Sigma} F(S)=\frac{18}{900}=\frac{1}{50},
\]
and the maximum is attained by exactly two configurations, namely
\[
 S_1=\{s:\ s\bmod 100\in\{7,93\}\},\qquad S_2=\{s:\ s\bmod 100\in\{43,57\}\},
\]
each consisting of $18$ types, all of weight $f(s,S_i)=1$. The second largest value is attained by
$S=\Z/900\Z$ (all types present) and equals
\[
 \frac{72K_2'+828J'}{900}=\frac{3.60820+11.79486}{900}=\frac{15.40306}{900}\approx 0.0171145<\frac{1}{50};
\]
this is also the largest value among configurations containing at least one type outside $C$.
\end{theorem}

The optimum was computed in three independent ways:
\begin{enumerate}
\item a CBC $0$--$1$ integer program on the reduced model ($72$ binary variables for $C$-types, $36$ for
the non-$C$ signature groups, plus full-weight indicator variables), by the author
(\texttt{typebound.py});
\item an exhaustive enumeration over the $64$ Galois-closed signature sets ($439$ valid pairs), with
each candidate realized as an explicit type set and $F$ evaluated directly from its definition on all
$900$ types (\texttt{opt\_closure.py}, referee);
\item a HiGHS MILP on the reduced model with no-good cuts, enumerating \emph{all} optimal solutions and
the next value below them (\texttt{opt\_milp\_reduced.py}, referee).
\end{enumerate}
No brute-force search over subsets of $\Z/900\Z$ was performed, and none is needed. As a
cross-validation, $3000$ random admissible configurations $S$ were evaluated directly and none
exceeded $18/900$.

\begin{remark}[The binding inequality and its robustness]\label{rem:binding}
Theorem~\ref{thm:opt} reduces the whole stability statement to the single numerical inequality
\[
 72K_2'+828J'\ <\ 18 ,
\]
which holds with room to spare: $72K_2'+828J'=15.40306$. At the true value of $K_2'$ the inequality
would still hold for any $J'<0.01738$. It is, however, genuinely necessary to use the refined constant
$J'$ rather than $K_1$: replacing $J'$ by $K_1$ gives $(72K_2'+828K_1)/900\approx 0.0292>1/50$.
\end{remark}

\section{Proof of Theorem~\ref{thm:main}}\label{sec:proof}

\begin{proof}[Proof of Theorem~\ref{thm:main}]
Let $\varepsilon>0$. Let $A\subseteq[N]$ be admissible, contained in neither extremal class, and let
$S\subseteq\Z/900\Z$ be its set of types. If some $a\in A$ has $a\not\equiv 7\ (25)$ then its type
$s$ satisfies $s\not\equiv 7\ (25)$, and similarly for $18$; hence $S\in\Sigma$.

By Proposition~\ref{prop:pertype} (applied with $\varepsilon/2$) and Theorem~\ref{thm:opt},
\[
 |A|\ \le\ N\,F(S)+\frac{\varepsilon}{2}N+O(900)\ \le\ \frac{N}{50}+\frac{\varepsilon}{2}N+O(900),
\]
which is at most $(1/50+\varepsilon)N$ once $N\ge N_0(\varepsilon)$ with $N_0$ large enough to absorb
the $O(900)$ term. This proves the upper bound.

Sharpness is Proposition~\ref{prop:extremal}: $A_1$ and $A_2$ are admissible, meet both classes, and
have $N/50+O(1)$ elements. That these are, within the mod-$900$ type framework, the only configurations
attaining the bound is the uniqueness statement of Theorem~\ref{thm:opt}, together with
Proposition~\ref{prop:only-two}.
\end{proof}

\section{Numerical evidence for small \texorpdfstring{$N$}{N}}\label{sec:small}

Exact maximum-clique computations in the compatibility graph on
$\{n\le N:\ n^2+1\ \text{not squarefree}\}$ give the following (these are consistent with, but of
course weak evidence for, an asymptotic statement):
\begin{itemize}
\item For every $N\le 10000$, the overall maximum admissible size equals
$|\{n\le N:\ n\equiv 7\ (25)\}|$, in agreement with Theorem~\ref{thm:li}.
\item For every $N$ with $557\le N\le 10000$, the largest admissible set not contained in a single
mod-$25$ class is exactly the $\{7,93\}$ (or $\{43,57\}$) mod $100$ set. For some $N\le 556$ it is one
element larger.
\item The largest admissible set containing an element whose type lies outside $C$ has about $0.0155N$
elements for $N=4000$--$10000$, below the asymptotic runner-up value $0.0171145$ of
Theorem~\ref{thm:opt}.
\end{itemize}

\section{Computational components and verification status}\label{sec:verif}

We separate clearly what is computed, what is verified, and what is taken as an input.

\paragraph{(A) Finite computations entering the proof of Theorem~\ref{thm:main}.}
\begin{enumerate}
\item \emph{Local and gluing checks mod $900$.} The facts (F1)--(F3) of Lemma~\ref{lem:F1}, the
signature reduction data (that all $23$ non-$C$ types of a signature behave identically towards each
$C$-type, and that the relation $R$ reproduces gluing in all cross-class and $C$/non-$C$ cases), and
Proposition~\ref{prop:only-two}, were verified by direct enumeration: (F1) for all primes $p<2000$,
(F2) on $20{,}000$ random pairs, (F3) and the signature data on all $900\times 900$ type pairs, and
the mod-$100$ cross-class pairs exhaustively (and against all integer pairs up to $20000$).
\item \emph{The mod-$900$ maximization.} Theorem~\ref{thm:opt} is a finite computation, performed by
the three independent solvers listed in Section~\ref{sec:opt} (CBC $0$--$1$ program; exhaustive
enumeration over the $64$ Galois-closed candidate configurations; and an independent HiGHS MILP solve). The three solvers were written independently and share
no code; they agree on the optimal value, on the extremal configuration, and on the certificate of
optimality (a dual/feasibility certificate in the $0$--$1$ program, and exhaustion in the other two).
% TODO(author): state explicitly which $6$-element Galois orbit structure the $64 = 2^6$ Galois-closed
% configurations enumerate, so that the second solver's search space is reproducible from the text alone.
\end{enumerate}

\paragraph{(B) The verified numerical range.}
The explicit certificate described in Section~\ref{sec:small} was generated and then re-checked by a
separate program. The generator emits, for each plateau of the relevant counting function, the
plateau endpoints and the associated witness data; the checker re-reads this output and verifies each
plateau independently of the generator's internal state, so that an error in the generator cannot be
absorbed silently. The verification covers all
\[
N \le 9{,}999{,}981,
\]
comprising $399{,}999$ plateaus. The resulting artefacts have MD5 checksums
\[
\texttt{e5b8573a}\ \ (\text{global certificate}),\qquad
\texttt{e3330948}\ \ (\text{plateau deltas}),
\]
quoted here by their leading bytes; the full checksums and both programs are distributed with the
code (see (D) below). We stress that this computation does not enter the proof of
Theorem~\ref{thm:main}: it is a consistency check on the constants and on the structure of the
extremal configuration in the accessible range.

\paragraph{(C) Independent rebuild of Li's formalization.}
The input Theorem~\ref{thm:li} relies on the Lean formalization of Li
\cite{Li2026,LiRepo}. We did not take the published build artefacts on trust. Instead we rebuilt the
project from source and replayed the kernel checks. Three points are worth recording.

\begin{enumerate}
\item \emph{Patch.} The repository as published did not compile against the toolchain we used; a
two declarations were missing the implicit binders they use, which the project's
\texttt{autoImplicit = false} setting turns into a build stall; a two-line patch adds the binders
$\{v\ w : \mathbb{N}\}$ and $\{w : \mathbb{N}\}$, with no change to any statement (the elaborated
types are identical, checked with \texttt{pp.all}). The patch is included verbatim in our repository so that the discrepancy between the
published source and the compiling source is visible rather than hidden.
\item \emph{Build.} With this patch the full project builds cleanly: $39{,}069$ jobs, exit code $0$,
no \texttt{sorry} and no errors.
\item \emph{Kernel replay and axiom audit.} We ran \verb|#print axioms| on the top-level statement
\verb|all_N| and obtained exactly
\[
\texttt{[propext, Classical.choice, Quot.sound]},
\]
i.e.\ the three standard axioms of Lean's \texttt{Mathlib} foundation and nothing further --- in
particular no \texttt{sorryAx} and no unverified native-code evaluation axiom. Independently, a run
of \texttt{leanchecker} over the compiled environment re-checked $30{,}636$ modules with $0$
failures, so the proof terms were replayed by the kernel rather than merely accepted from the
compiler's cache.
\end{enumerate}

Thus the input we take from \cite{Li2026,LiRepo} is a machine-checked theorem whose dependency set we
have audited, not an informal assertion.

\paragraph{(D) Availability.}
All code described above --- the three mod-$900$ solvers, the enumeration scripts for (F1)--(F3) and
for the signature reduction, the certificate generator and the independent certificate checker, the
build and kernel-replay scripts for (C), and the binder patch --- is publicly available at
\begin{center}
\url{https://github.com/Paulee314/erdos848}
\end{center}
together with the raw output files whose checksums are quoted in (B).

\paragraph{Summary of logical status.}
Theorem~\ref{thm:main} depends on: (i) the finite checks (A), each of which is an exhaustive
enumeration over an explicitly bounded search space; (ii) the formalized result of Li
\cite{Li2026,LiRepo}, audited as in (C); and (iii) the stability statement of Sawhney
\cite{Sawhney}. It does \emph{not} depend on the numerical range (B), which is reported as
independent corroborating evidence.

\begin{thebibliography}{9}

\bibitem{ErdosProblem848}
\emph{Erd\H{o}s Problem 848}, Erd\H{o}s Problems database.
\url{https://www.erdosproblems.com/848}.

\bibitem{Li2026}
A.~C.~Li,
\emph{Erd\H{o}s Problem 848: a kernel-checked proof of the exact extremal bound},
SSRN preprint, abstract id.~7230480.
\url{https://papers.ssrn.com/sol3/papers.cfm?abstract_id=7230480}.

\bibitem{LiRepo}
A.~C.~Li,
\emph{erdos-848-squarefree-product} (Lean formalization, source repository).
\url{https://github.com/crabsatellite/erdos-848-squarefree-product}.

\bibitem{Sawhney}
M.~Sawhney,
\emph{Stability for Erd\H{o}s Problem 848}.
As reported in the project sources accompanying \cite{ErdosProblem848}.
% TODO(author): replace with full bibliographic data (venue/preprint number/date) once the
% primary source for Sawhney's stability result is located; the project sources used here
% report the statement but give no further bibliographic details.

\bibitem{AstroCat}
AstroCat93,
\emph{erdos-848-verification} (GitHub repository).
\url{https://github.com/AstroCat93/erdos-848-verification}.

\bibitem{LeeRepo}
P.~J.~Lee,
\emph{erdos848}: solvers, certificates and verification scripts, 2026.
\url{https://github.com/Paulee314/erdos848}.

\end{thebibliography}

\end{document}